\chapter{研究内容及拟解决的关键问题}

\section{总体目标}
\placeholder{用一段话概括总体目标、研究对象、数据尺度、最终输出与应用场景。}

\section{研究内容}

\subsection{任务一：低空逐株出苗检测与计数}
\placeholder{以 12、15、20 m 等低空影像为主，研究小目标玉米幼苗的检测、定位与数量统计，并分析高度变化对目标可辨识性与计数精度的影响。}

\subsection{任务二：中高空作物行与断行/缺苗区域检测}
\placeholder{以 25--50 m 影像为主，从逐株检测转向行结构与缺苗片段识别，实现大区域快速筛查和空间定位。}

\subsection{任务三：单株冠层实例分割与影像表型测量}
\placeholder{对低空可分辨单株进行实例分割，提取冠层面积、形状、覆盖度等稳定的影像表型指标，避免无法由二维影像可靠反演的指标。}

\subsection{任务四：春玉米—夏玉米等跨域泛化与适用性验证}
\placeholder{构建跨区域或跨种植制度验证，研究模型在背景、土壤、光照、品种和栽培方式变化下的性能退化与泛化能力。}

\begin{table}[htbp]
  \centering
  \caption{四项研究任务的逻辑关系与主要输出（模板）}
  \begin{tabular}{p{0.11\textwidth}p{0.19\textwidth}p{0.25\textwidth}p{0.29\textwidth}}
    \toprule
    任务 & 主要尺度 & 核心问题 & 主要输出 \\
    \midrule
    任务一 & 12--20 m & 逐株可见条件下的小目标检测与计数 & 单株位置、出苗数、计数误差 \\
    任务二 & 25--50 m & 单株退化条件下的行结构与断行识别 & 作物行、断行、缺苗区域 \\
    任务三 & 低空单株 & 冠层边界与稳定影像表型 & 实例掩膜、冠层形态指标 \\
    任务四 & 跨区域/跨季节 & 分布变化下的模型稳健性 & 泛化性能、适用边界与迁移策略 \\
    \bottomrule
  \end{tabular}
\end{table}

\section{拟解决的关键问题}
\begin{enumerate}[label=（\arabic*）,leftmargin=3.2em]
  \item \placeholder{多飞行高度导致目标尺度和纹理信息显著变化时，如何保持低空逐株检测的稳定性。}
  \item \placeholder{高空影像中单株目标弱化后，如何利用行结构与空间连续性实现缺苗/断行识别。}
  \item \placeholder{如何界定并提取对分辨率、遮挡和光照变化较稳健的单株影像表型。}
  \item \placeholder{如何评估并提升模型在春玉米—夏玉米、不同区域和不同栽培背景下的跨域泛化能力。}
\end{enumerate}

\section{四项任务之间的逻辑关系}
\placeholder{建议用“低空精细感知提供个体基准—中高空结构识别扩大覆盖—实例表型补充苗情质量维度—跨域验证形成可推广方法”的链条说明。}
